\begin{abstract}
Mathematical reasoning remains a formidable task for language models, primarily due to its intricate and highly structured requirements. In this work, we present DeepSeekMath~7B, an extension of DeepSeek-Coder-Base-v1.5 7B, which undergoes additional pre-training on 120B math-focused tokens compiled from Common Crawl, supplemented by both natural language and code data. Without employing external auxiliary tools or ensemble techniques, DeepSeekMath~7B attains a remarkable 51.7\% score on the MATH competition-level benchmark, nearing the capabilities of top-tier models like Gemini-Ultra and GPT-4. When self-consistency is applied over 64 samples, DeepSeekMath~7B reaches 60.9\% on MATH. The advanced mathematical reasoning skills demonstrated by DeepSeekMath~stem from two principal strategies: (1) leveraging large-scale public web datasets via a carefully constructed data selection process, and (2) implementing Group Relative Policy Optimization (GRPO)—an adaptation of Proximal Policy Optimization (PPO)—which boosts mathematical reasoning efficiency while improving PPO’s memory performance.
\end{abstract}

\section{Introduction}

Large language models (LLMs) have transformed artificial intelligence approaches to mathematical reasoning, driving substantial progress on quantitative reasoning tasks \citep{MATH} and geometry-based evaluations \citep{trinh2024solving}. These systems are increasingly valuable for supporting human users in addressing advanced mathematical challenges \citep{tao}. Nevertheless, leading-edge models such as GPT-4 \citep{gpt4} and Gemini-Ultra \citep{gemini} are not openly available to the research community, and the performance of available open-source models is markedly inferior in comparison.

This work presents DeepSeekMath, a specialized language model tailored for mathematical domains that demonstrates substantial improvements over open-source alternatives and attains near-GPT-4 performance on a range of academic benchmarks. Central to our approach is the construction of the DeepSeekMath Corpus, an expansive, high-quality pre-training dataset encompassing 120B mathematical tokens. To assemble this corpus, we extract data from the Common Crawl (CC) by leveraging a classifier grounded in fastText \citep{joulin2016fasttext}. During its first phase, the classifier is trained on positive samples from OpenWebMath \citep{openwebmath}, combined with a heterogeneous assortment of other web pages as negative instances. After this initial stage, the classifier is utilized to identify additional candidate positives from CC; these selections are subsequently curated through manual annotation. Incorporating this improved data, the classifier is iteratively refined. Our evaluation demonstrates that this large-scale corpus is of outstanding quality: DeepSeekMath-Base 7B attains 64.2\% on GSM8K \citep{gsm8k} and 36.2\% on the high-difficulty MATH benchmark \citep{MATH}, surpassing Minerva 540B \citep{minerva}. Furthermore, due to the multilingual nature of the DeepSeekMath Corpus, we also observe enhanced results on Chinese mathematical datasets \citep{wei2023cmath,agieval}. We suggest that the methodologies explored here for mathematical data processing offer a valuable reference for the broader research community and foresee considerable opportunities for further advancement.

DeepSeekMath-Base is constructed by initializing with DeepSeek-Coder-Base-v1.5 7B \citep{deepseek-coder}, given our observation that a foundation in code-oriented pretraining yields superior results compared to models derived from generic large language models. Additionally, we observe that subsequent mathematical training not only sharpens mathematical proficiency but also leads to better performance on benchmarks such as MMLU \citep{mmlu} and BBH \citep{bbh}, signifying a broader improvement in reasoning abilities beyond mathematics.

Following the pre-training phase, we further fine-tune DeepSeekMath-Base using mathematical instruction data encompassing chain-of-thought \citep{cot}, program-of-thought \citep{pot,pal}, and tool-integrated reasoning \citep{tora} approaches. This process produces DeepSeekMath-Instruct 7B, which surpasses all other 7B-sized models and approaches the capability of the top-tier 70B open-source instruction-tuned models.

In addition, we propose Group Relative Policy Optimization (GRPO), a new reinforcement learning (RL) algorithm related to Proximal Policy Optimization (PPO) \citep{schulman2017proximal}. GRPO eliminates the necessity for a separate critic network and, instead, computes the baseline directly from group scores, enabling significant resource efficiency during training. Utilizing only a portion of English instruction tuning data, GRPO considerably advances the already strong DeepSeekMath-Instruct’s performance—showing marked improvements on in-domain tasks (GSM8K: 82.9\% $\rightarrow$ 88.2\%, MATH: 46.8\% $\rightarrow$ 51.7\%) and out-of-domain settings (e.g., CMATH: 84.6\% $\rightarrow$ 88.8\%) throughout the reinforcement learning process. We further establish a cohesive framework for interpreting several approaches—including Rejection Sampling Fine-Tuning (RFT) \citep{yuan2023scaling}, Direct Preference Optimization (DPO) \citep{dpo}, PPO, and GRPO—revealing these methods as variations of direct or simplified RL formulations. Through extensive experimental analysis—addressing comparisons such as online versus offline learning, outcome-based versus process-based supervision, and single-turn versus iterative reinforcement learning—we thoroughly explore the core factors underpinning this framework. Finally, we elucidate the mechanisms by which our RL paradigm enhances instruction-tuned model performance and suggest future avenues for realizing more efficient RL methodologies within this unified conceptual structure.

\subsection{Contributions}
This work presents advancements in large-scale mathematical pre-training and a systematic investigation of reinforcement learning strategies.

\textbf{Math Pre-Training at Scale}

\begin{itemize}[topsep=0pt]
    \item We demonstrate that the open-access Common Crawl repository is a rich resource for mathematical content. Through the development of a comprehensive data selection pipeline, we curate the DeepSeekMath Corpus, a robust dataset comprising 120B tokens drawn from web pages with significant mathematical relevance. This corpus substantially surpasses prior datasets in scale, being nearly sevenfold larger than the math-specific web data of Minerva \citep{minerva} and approximately nine times greater than that in the OpenWebMath release \citep{openwebmath}.

    \item Our pre-trained DeepSeekMath-Base 7B model matches the mathematical reasoning performance of the much larger Minerva 540B \citep{minerva}. This finding suggests that sheer model scale is not solely responsible for superior performance in mathematical domains; in fact, a well-curated, high-quality dataset allows smaller models to achieve competitive results.

    \item We report results from various mathematical pre-training strategies. Pre-training on code before proceeding to math-specific data notably enhances the model's problem-solving abilities, both in pure reasoning tasks and when tools are involved. This addresses, in part, the enduring debate: \textit{does code pre-training aid reasoning skills?} For mathematical reasoning, our evidence indicates a positive impact.

    \item While many previous works incorporate arXiv papers into the training corpus, our analysis shows that such data does not yield significant improvements across any of the mathematical benchmarks considered here.
\end{itemize}

\textbf{Exploration and Analysis of Reinforcement Learning}

\begin{itemize}[topsep=0pt]
    \item We present Group Relative Policy Optimization (GRPO), a reinforcement learning algorithm designed for both efficiency and efficacy. Distinct from methods relying on a critic network, GRPO computes the baseline using aggregated group scores, thereby requiring substantially fewer training resources than Proximal Policy Optimization (PPO).
    \item Through empirical analysis, we show that GRPO markedly improves the performance of our instruction-tuned model DeepSeekMath-Instruct using solely instruction-tuning data. Additionally, we identify gains in generalization to out-of-domain tasks throughout reinforcement learning.
    \item We propose a cohesive framework to analyze and relate various approaches, including RFT, DPO, PPO, and GRPO. Comprehensive experimentation is conducted—contrasting online versus offline learning, evaluating outcome-based against process-based supervision, and examining both single-turn and iterative reinforcement learning—to rigorously dissect the underlying mechanics of the framework.
    \item Leveraging insights from this unified perspective, we investigate the mechanisms driving the success of reinforcement learning, outlining prospective strategies for further enhancing reinforcement learning with large language models (LLMs).
\end{itemize}

\subsection{Summary of Evaluations and Metrics}

\begin{itemize}[topsep=0pt]
    \item \textbf{English and Chinese Mathematical Reasoning}: Our evaluation encompasses a broad spectrum of mathematical reasoning benchmarks in both English and Chinese, ranging from elementary to undergraduate-level problems. English benchmarks cover datasets such as GSM8K \citep{gsm8k}, MATH \citep{MATH}, SAT \citep{llemma}, OCW Courses \citep{minerva}, and MMLU-STEM \citep{mmlu}. The Chinese benchmarks used include MGSM-zh \citep{mgsm}, CMATH \citep{wei2023cmath}, Gaokao-MathCloze \citep{agieval}, and Gaokao-MathQA \citep{agieval}. We measure models' proficiency both in producing fully self-contained written solutions and in addressing problems computationally via Python.

    On the English datasets, DeepSeekMath-Base achieves performance on par with the proprietary Minerva 540B \citep{minerva}, and significantly outperforms all open-source base models (for example, Mistral 7B \citep{mistral} and Llemma-34B \citep{llemma}), regardless of the presence of math-specific pre-training, frequently by a considerable gap. Significantly, DeepSeekMath-Base leads on the Chinese benchmarks, attributable to our deviation from prior work \citep{minerva,llemma} by incorporating high-quality math pre-training data from multiple languages rather than restricting to English. When enhanced with mathematical instruction tuning and reinforcement learning, the resultant DeepSeekMath-Instruct and DeepSeekMath-RL achieve outstanding results, for the first time enabling an open-source model to surpass 50\% accuracy on the challenging, competition-grade MATH dataset.
\end{itemize}

\item \textbf{Formal Mathematics}: We assess DeepSeekMath-Base's capability in informal-to-formal theorem proving by utilizing the task from \citep{dsp_proof} on the miniF2F benchmark \citep{minif2f}, with Isabelle \citep{isabelle} serving as the proof assistant. DeepSeekMath-Base achieves notable results in few-shot autoformalization scenarios.

\item \textbf{Natural Language Understanding, Reasoning, and Code}: To provide a thorough evaluation of the model's abilities in comprehension, logical reasoning, and programming, we examine DeepSeekMath-Base's performance on several benchmarks. These include the Massive Multitask Language Understanding (MMLU) dataset \citep{mmlu}, which spans 57 multiple-choice tasks across a range of topics, and BIG-Bench Hard (BBH) \citep{bbh}, which focuses on 23 tasks designed to test complex, multi-step reasoning. For code-related capabilities, we use HumanEval \citep{codex} and MBPP \citep{mbpp}, both established for the assessment of code generation models. The results demonstrate that math-oriented pre-training enhances both linguistic understanding and reasoning skills.

\section{Math Pre-Training}

\subsection{Data Collection and Decontamination}
Here, we describe the methodology for constructing the DeepSeekMath Corpus using data sourced from Common Crawl. As illustrated in Figure \ref{fig:data_collect}, we employ an iterative data pipeline that enables the large-scale extraction of mathematical content from Common Crawl, initiating the process with a curated, high-quality seed set of mathematical texts. This methodology can similarly be applied to domains such as programming.

Our initial step involves selecting OpenWebMath \citep{openwebmath}—a vetted collection of mathematical text from the web—as the seed dataset. With this collection, we train a fastText model \citep{joulin2016fasttext} to facilitate retrieval of additional mathematical web pages with similar characteristics. Specifically, we sample 500,000 documents from OpenWebMath as positive examples and pair these with 500,000 negative samples drawn randomly from Common Crawl. Model training is conducted using an open-source implementation\footnote{\url{https://fasttext.cc}}, with settings including a 256-dimensional vector space, a learning rate of 0.1, a maximum word n-gram length of 3, a minimum occurrence threshold of 3, and 3 epochs. 

To further trim Common Crawl's size, we implement both exact and near-duplicate URL removal strategies, producing a dataset of 40B HTML pages. Utilizing the trained fastText classifier, we retrieve mathematical web pages from this deduplicated corpus. We then rank these pages by the scores generated from the fastText model, retaining only those that score highest in mathematical relevance. The quantity of data selected for pre-training is empirically determined via experiments using the top 40B, 80B, 120B, and 160B tokens, with the initial iteration retaining the top 40B tokens.

Following the initial round of data collection, a significant number of mathematical web pages remain uncollected. This is largely attributed to the limited diversity of positive examples used to train the fastText model. To address this, we expand the seed corpus by sourcing additional mathematical websites, thereby enhancing the fastText model’s effectiveness. Our procedure involves segmenting the full Common Crawl into non-overlapping domains, where each domain groups web pages by their base URL. For every domain, we determine what proportion of its pages were collected during the first round. Those domains in which more than 10\% of pages were gathered are categorized as being math-focused (for instance, \url{mathoverflow.net}). 

We then proceed to manually label URLs within these selected domains that are known to feature mathematical content (such as \url{mathoverflow.net/questions}). Web pages associated with these identified URLs but previously uncollected are subsequently incorporated into the seed corpus. This method allows us to broaden our set of positive samples, supporting the training of a more robust fastText model capable of retrieving additional mathematical web data in further iterations. Ultimately, after four cycles of data gathering, we compile a dataset consisting of 35.5 million mathematical web pages and a cumulative total of 120 billion tokens. Notably, by the fourth collection iteration, almost 98\% of the corpus was already assembled in the third iteration, leading us to discontinue further collection.

To safeguard against benchmark data leakage, we employ the filtering strategy described in \cite{deepseek-coder}, which involves excluding web pages that include questions or answers present in English mathematical benchmarks such as GSM8K \citep{gsm8k} and MATH \citep{MATH}, as well as Chinese benchmarks including CMATH \citep{wei2023cmath} and AGIEval \citep{agieval}. Specifically, any textual passage that features a 10-gram string with an exact match to any substring from these evaluation benchmarks is purged from our training corpus. For benchmark sequences shorter than 10 grams but at least 3 grams in length, we implement precise matching to exclude potentially contaminated content.

\subsection{Validating the Quality of the DeepSeekMath~Corpus}

To assess the quality of DeepSeekMath~Corpus, we conduct pre-training studies comparing it to the latest mathematical corpora:

\begin{itemize}[topsep=0pt]
    \item \textbf{MathPile} \citep{mathpile}: a corpus combining multiple sources (totaling 8.9B tokens) from textbooks, Wikipedia, ProofWiki, CommonCrawl, StackExchange, and arXiv, with arXiv comprising more than 85\% of the tokens;
    \item \textbf{OpenWebMath} \citep{openwebmath}: a dataset of mathematical content extracted from CommonCrawl, containing 13.6B tokens;
    \item \textbf{Proof-Pile-2} \citep{llemma}: a math dataset comprising OpenWebMath, AlgebraicStack (10.3B tokens of math code), and arXiv documents (28.0B tokens). For Proof-Pile-2 experiments, we adopt the arXiv:Web:Code distribution of 2:4:1, as indicated in \cite{llemma}.
\end{itemize}

\subsubsection{Training Setting}
\label{sec:quality-policy}
We conduct mathematical training on a generic pre-trained language model with 1.3B parameters, referred to as DeepSeek-LLM 1.3B, which utilizes the same architecture as DeepSeek LLMs \citep{deepseek-llm}. For each mathematical dataset, a separate model is trained for 150B tokens. All model training is performed using the HAI-LLM framework \citep{haillm}, chosen for its efficiency and lightweight nature. Consistent with DeepSeek LLM training protocols, we adopt the AdamW optimizer \citep{adamW} with $\beta_1=0.9$, $\beta_2=0.95$, and $\mathrm{weight\_decay}=0.1$. A multi-step learning rate policy is implemented: the learning rate peaks after 2,000 warmup steps, reduces to 31.6\% of its maximum at 80\% of total training steps, and further declines to 10.0\% by 90\% of training completion. The maximum learning rate is set at 5.3e-4, with training performed using a batch size of 4M tokens and a sequence context window of 4K tokens.

\subsubsection{Evaluation Results}

\textbf{The DeepSeekMath~Corpus excels in quality, supports multiple languages in mathematical content, and is distinguished by its large scale.}

\begin{itemize}[topsep=0pt]
    \item \textbf{High-quality}: 
    We assess downstream task performance on eight mathematical benchmarks via few-shot chain-of-thought prompting \cite{cot}. Results displayed in Table \ref{tab:corpora_comparison} illustrate a substantial performance advantage for the model trained on DeepSeekMath~Corpus. As depicted in Figure \ref{fig:corpora_comparisons}, the model using DeepSeekMath Corpus outperforms the model trained on Proof-Pile-2 after 50B tokens (equivalent to a single epoch on Proof-Pile-2), suggesting that DeepSeekMath Corpus maintains a superior average quality.
    \item \textbf{Multilingual}: 
    The DeepSeekMath~Corpus consists of data in a variety of languages, with English and Chinese forming the bulk of the dataset. Table \ref{tab:corpora_comparison} demonstrates that training on the DeepSeekMath~Corpus leads to notable gains in mathematical reasoning for both English and Chinese. On the other hand, alternative mathematical corpora, being predominantly English-focused, contribute less to improvements and may even restrict progress in Chinese mathematical reasoning tasks.
    \item \textbf{Large-scale}: 
    The scale of DeepSeekMath~Corpus greatly exceeds that of any current mathematical corpus. As indicated in Figure \ref{fig:corpora_comparisons}, models trained on DeepSeekMath~Corpus, specifically DeepSeek-LLM 1.3B, exhibit a faster and more sustained learning progression. In contrast, the baseline datasets, which are much smaller in size and undergo multiple repeated passes during training, quickly yield models whose performance plateaus.
\end{itemize}

\subsection{Training and Evaluating DeepSeekMath-Base 7B}

This section presents DeepSeekMath-Base 7B, a foundational model designed to exhibit robust mathematical reasoning. The model builds on the weights of DeepSeek-Coder-Base-v1.5 7B \citep{deepseek-coder}, serving as its initialization point, and is further trained over a total of 500B tokens. The training data is comprised of 56\% DeepSeekMath Corpus, 4\% AlgebraicStack, 10\% arXiv articles, 20\% Github code, and 10\% natural language content sourced from the Common Crawl dataset, incorporating both English and Chinese texts. The training approach largely mirrors that outlined in Section \ref{sec:quality-policy}, with the exceptions that the learning rate peaks at 4.2e-4, and training is performed with a batch size of 10M tokens.

A thorough evaluation is performed to ascertain DeepSeekMath-Base 7B's mathematical proficiency. This assessment targets its capability to independently generate comprehensive mathematical solutions without tool assistance, address mathematical questions utilizing computational tools, and execute formal theorem proving. Furthermore, we broaden our analysis to survey the base model’s performance in areas such as natural language understanding, reasoning aptitude, and programming competence.

\paragraph{Mathematical Problem Solving with Step-by-Step Reasoning}

To determine the model’s aptitude for mathematical problem solving, we test DeepSeekMath-Base using few-shot chain-of-thought prompting \citep{cot} on eight English and Chinese benchmarks. These benchmarks test quantitative reasoning—using datasets such as GSM8K \citep{gsm8k}, MATH \citep{MATH}, and CMATH \citep{wei2023cmath}—as well as multiple-choice assessments like MMLU-STEM \citep{mmlu} and Gaokao-MathQA \citep{agieval}, encompassing a variety of mathematical topics ranging from elementary to university-level difficulty.

As indicated in Table \ref{tab:base_math_cot}, DeepSeekMath-Base 7B consistently achieves the highest scores across all eight assessed benchmarks among open-source base models, which include both the popular Mistral 7B \citep{mistral} and the newly introduced Llemma 34B \citep{llemma}, the latter of which was specifically trained on mathematical data from Proof-Pile-2 \citep{llemma}. Particularly striking is its performance on the MATH dataset, where DeepSeekMath-Base exceeds the accuracy of other available open-source models by more than 10 percentage points. Furthermore, it outperforms Minerva 540B \citep{minerva}—a closed-source model of much greater scale (77 times larger) constructed on PaLM \citep{palm} and subsequently refined using mathematical literature.

\paragraph{Mathematical Problem Solving with Tool Use} We assess the models' ability to perform program-assisted mathematical reasoning on GSM8K and MATH, employing few-shot program-of-thought prompting techniques \citep{pot,pal}. Here, the models are directed to generate Python scripts, drawing on computational packages such as \textit{math} and \textit{sympy} for more complex tasks, with the answer determined by the program’s output. As depicted in Table \ref{tab:base_math_pot_proof}, DeepSeekMath-Base 7B delivers superior results, surpassing the previous best, Llemma 34B.

\paragraph{Formal Mathematics}

Automating formal proof generation has become increasingly important for bolstering both the reliability and efficiency of mathematical verification. We measure DeepSeekMath-Base 7B's competence on informal-to-formal proof translation following the setup in \citep{dsp_proof}: the task requires producing a formal proof given an informal description, a formal version of the statement, and an informal argument. Evaluation is performed on miniF2F \citep{minif2f}, a benchmark for formalized Olympiad-level mathematics, where the model is asked to construct proofs in Isabelle using a few exemplars. Consistent with the methodology in \cite{dsp_proof}, the approach involves having the model generate an initial proof outline, after which the Sledgehammer automated theorem prover \citep{sledgehammer} fills in any details left unspecified. The outcomes summarized in Table \ref{tab:base_math_pot_proof} demonstrate that DeepSeekMath-Base 7B excels in the autoformalization of mathematical proofs.

\paragraph{Natural Language Understanding, Reasoning, and Code}

Model effectiveness in natural language comprehension is benchmarked on MMLU \citep{mmlu}, reasoning skills on BBH \citep{bbh}, and programming capabilities on HumanEval \citep{codex} as well as MBPP \citep{mbpp}. Results presented in Table \ref{tab:understanding_reasoning_code} reveal that DeepSeekMath-Base 7B achieves notable gains on MMLU and BBH compared to its predecessor, DeepSeek-Coder-Base-v1.5 \citep{deepseek-coder}, highlighting how mathematical training enhances both language and reasoning proficiencies. Moreover, the integration of code-related tokens during continual training allows DeepSeekMath-Base 7B to sustain the same coding benchmark results as DeepSeek-Coder-Base-v1.5. In all, DeepSeekMath-Base 7B outshines the general-purpose Mistral 7B \citep{mistral} by a substantial margin across the reasoning and coding tasks considered.

\section{Supervised Fine-Tuning}

\subsection{SFT Data Curation}

We assemble an instruction-tuning dataset for mathematics that encompasses a broad range of problems in both English and Chinese, drawing from multiple mathematical domains and difficulty levels. Each problem is matched with a corresponding solution presented in chain-of-thought (CoT) \citep{cot}, program-of-thought (PoT) \citep{pot,pal}, or tool-integrated reasoning format \citep{tora}. The compiled dataset comprises 776,000 training instances.

\begin{itemize}[topsep=0pt]
    \item \textbf{English mathematical datasets}: We provide tool-integrated solution annotations for GSM8K and MATH datasets and supplement these with a selected portion of MathInstruct \citep{MathInstruct}, alongside the training split from Lila-OOD \citep{lila}, where answers are constructed via CoT or PoT strategies. The English subset incorporates a wide array of mathematical areas, such as algebra, probability, number theory, calculus, and geometry.
    \item \textbf{Chinese mathematical datasets}: For Chinese, we gather K-12 mathematical questions spanning 76 distinct sub-topics (e.g., linear equations) and supply solutions employing both CoT and tool-integrated reasoning.
\end{itemize}

\subsection{Training and Evaluating DeepSeekMath-Instruct 7B}

This section details the supervised instruction tuning process applied to DeepSeekMath-Instruct 7B, built upon DeepSeekMath-Base. Training samples are concatenated at random until the context window reaches 4K tokens. The fine-tuning process utilizes a constant learning rate of 5e-5, a batch size of 256, and proceeds for 500 optimization steps.

Model evaluation encompasses both tool-augmented and standalone mathematical reasoning, using four quantitative reasoning benchmarks in English and Chinese. We position our model in comparison to other state-of-the-art systems of the period:

\begin{itemize}[topsep=0pt]
    \item \textbf{Closed-source models}: The evaluation set comprises: (1) models from the GPT series, notably GPT-4 \citep{gpt4} and GPT-4 Code Interpreter~\footnote{\url{https://openai.com/blog/chatgpt-plugins\#\#code-interpreter}}; (2) Gemini Ultra and Pro \citep{gemini}; (3) Inflection-2 \citep{inflection-2}; (4) Grok-1~\footnote{\url{https://x.ai/model-card}}; in addition to models recently launched by Chinese technology companies such as (5) Baichuan-3~\footnote{\url{https://www.baichuan-ai.com}} and (6) the most recent GLM-4~\footnote{\url{https://open.bigmodel.cn/dev/api\#glm-4}}

    from the GLM model family \citep{glm}.

These models are designed for broad usage, and most have undergone extensive alignment processes. 
\item \textbf{Open-source models} encompass: general-purpose models such as (1) DeepSeek-LLM-Chat 67B \citep{deepseek-llm}, (2) Qwen 72B \citep{qwen}, (3) SeaLLM-v2 7B \citep{seallm}, and (4) ChatGLM3 6B \citep{chatglm3}, alongside models tailored for mathematical reasoning. This latter group includes (5) InternLM2-Math 20B~\footnote{\url{https://github.com/InternLM/InternLM-Math}}, which extends InternLM2 through additional mathematics pretraining and instruction fine-tuning, (6) Math-Shepherd-Mistral 7B, which employs PPO optimization \citep{schulman2017proximal} on Mistral 7B \citep{mistral} with a supervision-by-process reward approach, (7) the WizardMath series \citep{wizardmath} that improves math reasoning abilities in Mistral 7B and Llama-2 70B \citep{llama2} via evolve-instruct—a variant of instruction tuning using machine-generated instructions—and PPO on datasets largely drawn from GSM8K and MATH, (8) MetaMath 70B \citep{metamath}, which refines Llama-2 70B with an expanded set from GSM8K and MATH, (9) ToRA 34B \cite{tora}, which adapts CodeLlama 34B for math reasoning augmented by tool usage, and (10) MAmmoTH 70B \citep{MathInstruct}, an instruction-tuned Llama-2 70B trained on MathInstruct.

As summarized in Table \ref{tab:sft_rl_math}, when models are assessed under a setting that prohibits tool usage, DeepSeekMath-Instruct 7B shows exceptional capacity for detailed, stepwise reasoning. On the advanced MATH benchmark, our approach outperforms all currently available open-source models and exceeds most proprietary offerings (such as Inflection-2 and Gemini Pro) by a margin of at least 9 percentage points. This holds even in comparison to much larger models (e.g., Qwen 72B) or those specifically optimized with reinforcement learning for mathematical skills (e.g., WizardMath-v1.1 7B). DeepSeekMath-Instruct’s results on MATH approach those of leading Chinese proprietary systems, including GLM-4 and Baichuan-3, but still trail the highest performers, namely GPT-4 and Gemini Ultra.

In evaluation conditions where both natural language reasoning and code-based tools can be used to tackle mathematical problems, DeepSeekMath-Instruct 7B attains almost 60\% accuracy on the MATH benchmark, surpassing all other open-source alternatives. For other datasets, its results are on par with DeepSeek-LLM-Chat 67B, which, notably, is ten times larger and previously led the field.

\section{Reinforcement Learning}

\subsection{Group Relative Policy Optimization} Reinforcement learning (RL) has demonstrated notable success in augmenting the mathematical reasoning abilities of LLMs following the Supervised Fine-Tuning (SFT) phase \citep{wang2023math,wizardmath}. Here, we present our RL framework, Group Relative Policy Optimization (GRPO), which is both scalable and effective.

\subsubsection{From PPO to GRPO} Proximal Policy Optimization (PPO) \citep{schulman2017proximal}, a widely adopted actor-critic RL algorithm, forms the backbone of many RL fine-tuning pipelines for LLMs \citep{ouyang2022training}. PPO optimizes LLM policies via the following surrogate loss function:

\begin{equation}
\footnotesize
    \mathcal{J}_{PPO}(\theta) = \mathbb{E}{[q \sim P(Q), o \sim \pi_{\theta_{old}}(O|q)]} \frac{1}{|o|} \sum_{t=1}^{|o|} \min \left[ \frac{\pi_\theta(o_{t} | q, o_{<t})}{\pi_{\theta_{old}}(o_{t} | q, o_{<t})} A_{t}, \text{clip} \left( \frac{\pi_\theta(o_{t} | q, o_{<t})}{\pi_{\theta_{old}}(o_{t} | q, o_{<t})}, 1 - \varepsilon, 1 + \

Here, $\pi_{\theta}$ represents the current policy model, while $\pi_{\theta_{old}}$ denotes the previous policy. The variables $q$ and $o$ refer to the questions and outputs, respectively, which are sampled from the dataset of questions and from the old policy $\pi_{\theta_{old}}$. The parameter $\varepsilon$ is a clipping hyper-parameter introduced in PPO to promote training stability. The term $A_t$ denotes the advantage, computed using Generalized Advantage Estimation (GAE) \citep{gae}, which leverages the sequence of rewards $\{r_{\ge t}\}$ along with a learned value function $V_{\psi}$. Consequently, PPO requires training an additional value function in parallel with the policy. To prevent the reward model from being overexploited, a standard technique involves including a token-level KL divergence penalty with respect to a reference policy in the reward at every generation step \citep{ouyang2022training}, as expressed by:

\begin{equation}
    r_{t} = r_\varphi(q, o_{\le t}) - \beta \log\frac{\pi_{\theta}(o_{t}|q, o_{<t})}{\pi_{ref}(o_{t}|q, o_{<t})},
\label{eq:PPO-reward}
\end{equation}

where $r_\varphi$ indicates the reward model, $\pi_{ref}$ specifies the reference model—typically the initial SFT policy—and $\beta$ determines the weight of the KL term.

Since PPO’s value function often matches the policy model’s size, this incurs significant computational and memory costs. Furthermore, during RL training, the value function serves as a baseline in the advantage computation to reduce variance. In practice for LLMs, the reward model commonly assigns a reward to just the last token of the output, making the value function difficult to train accurately across all tokens. To tackle this, we introduce Group Relative Policy Optimization (GRPO), depicted in Figure \ref{fig:grpo}, which circumvents the need for an explicit value function. Instead, GRPO uses the mean reward of several sampled completions—each generated for the same prompt—as a baseline. Concretely, for every question $q$, GRPO samples a batch of outputs $\{o_1, o_2, \cdots, o_G\}$ using the old policy $\pi_{\theta_{old}}$, and then the policy is optimized by maximizing:

\begin{equation}
\footnotesize
\begin{split}
    \mathcal{J}_{GRPO}(\theta) &= \mathbb{E}{[q \sim P(Q), \{o_i\}_{i=1}^G \sim \pi_{\theta_{old}}(O|q)]}  \\
    & \frac{1}{G}\sum_{i=1}^G\frac{1}{|o_i|} \sum_{t=1}^{|o_i|} \left\{ \min \left[ \frac{\pi_\theta(o_{i,t} | q, o_{i,<t})}{\pi_{\theta_{old}}(o_{i,t} | q, o_{i,<t})} \hat{A}_{i,t}, \text{clip} \left( \frac{\pi_\theta(o_{i,t} | q, o_{i,<t})}{\pi_{\theta_{old}}(o_{i,t} | q, o_{i,<t})}, 1 - \varepsilon, 1 + \varepsilon \right)  \hat{A}_{i,t} \right] - \beta \mathbb{D}_{KL}\left[\pi_{\theta} || \pi_{ref}\right]\right\} ,
\end{split}
\label{eq:GRPO-obj}
\end{equation}

In this formulation, $\varepsilon$ and $\beta$ are again hyper-parameters, and $\hat{A}_{i,t}$ represents the advantage calculated solely from the relative rewards among the outputs in each group; additional details are provided in the subsequent sections. By evaluating advantages within groups, GRPO directly leverages the comparative setup on which most reward models are trained, since such models generally rely on side-by-side comparisons of multiple answers to the same question. Furthermore, rather than embedding the KL penalty into the reward, GRPO penalizes divergence by explicitly including the KL term between the learned policy and the reference policy in the loss function, thereby simplifying the calculation of $\hat{A}_{i,t}$. Importantly, unlike the KL penalty used in (\ref{eq:PPO-reward}), we estimate KL divergence via the following unbiased estimator \citep{kl_approx}:

\begin

\begin{algorithm}[t]
  \small
  \caption{Iterative Group Relative Policy Optimization}
  \textbf{Input} initial policy model $\pi_{\theta_{\text{init}}}$; reward models $r_\varphi$; task prompts $\mathcal{D}$; 
  hyperparameters $\varepsilon$, $\beta$, $\mu$
  \begin{algorithmic}[1]
    \State Initialize the policy model: $\pi_\theta \leftarrow \pi_{\theta_{\text{init}}}$
    \For{iteration = 1, \dots, I}
      \State Assign the reference model: $\pi_{ref} \leftarrow \pi_{\theta}$
      \For{step = 1, \dots, M}
      \State Draw a minibatch $\mathcal{D}_b$ from $\mathcal{D}$
      \State Set the old policy model: $\pi_{\theta_{old}} \leftarrow \pi_{\theta}$ 
      \State For each prompt $q$ in $\mathcal{D}_b$, sample $G$ outputs $\{o_i\}_{i=1}^G \sim \pi_{\theta_{old}} (\cdot \mid q) $
      \State For every generated output $o_i$, use $r_{\varphi}$ to obtain rewards $\{r_i\}_{i=1}^{G}$
      \State For each output $o_i$, estimate the group relative advantage $\hat{A}_{i,t}$ for the $t$-th token
      \For{GRPO iteration = 1, \dots, $\mu$}
        \State Update $\pi_{\theta}$ by maximizing the GRPO objective (Equation \ref{eq:GC-GRPO})
      \EndFor
    \EndFor 
    \State Periodically update the reward model $r_\varphi$ through continual training leveraging a replay strategy. 
    \EndFor 
  \end{algorithmic}
  \textbf{Output} $\pi_\theta$
  \label{alg:iter-grpo}
\end{algorithm}

\subsubsection{Outcome Supervision RL with GRPO} Consider, for a given question $q$, generating a collection of outputs $\{o_1, o_2, \cdots, o_G\}$ using the old policy model $\pi_{\theta_{old}}$. Each output is then assigned a reward through the reward model, resulting in $G$ rewards $\mathbf{r} = \{r_1, r_2, \cdots, r_G\}$. These rewards are subsequently standardized by removing their mean and scaling by their standard deviation. Under the outcome supervision regime, the standardized reward associated with each output $o_i$ is assigned to all of its tokens as their advantage: $\hat{A}_{i, t} = \widetilde{r}_i = \frac{r_i- {\rm mean}(\mathbf{r})}{{\rm std}(\mathbf{r})}$. The policy is then updated by maximizing the objective in equation (\ref{eq:GRPO-obj}).

\subsubsection{Process Supervision RL with GRPO} In outcome supervision, rewards are provided solely at the conclusion of each output, which may limit the capacity to effectively instruct the model during sophisticated mathematical problem-solving. As in \cite{wang2023math}, we extend this by examining process supervision, wherein a reward is furnished at each intermediate reasoning step. Given a question $q$ and the batch of $G$ generated outputs $\{o_1, o_2, \cdots, o_G\}$, a process reward model scores each step, yielding sets of step-level rewards: $\mathbf{R} = \{ \{r_1^{index(1)},\cdots,r_1^{index(K_1)}\}, \cdots,  \{r_G^{index(1)},\cdots,r_G^{index(K_G)}\} \}$, with $index(j)$ indicating the endpoint of the $j$-th reasoning step, and $K_i$ the number of steps in output $i$. The step rewards are normalized: $\widetilde{r}_i^{index(j)} = \frac{r_i^{index(j)} - {\rm mean(\mathbf{R})}}{{\rm std(\mathbf{R})}}$. Next, the advantage for each token is computed as the sum of the normalized rewards from subsequent steps: $\hat{A}_{i, t} = \sum_{index(j) \ge t} \widetilde{r}_i^{index(j)}$. Policy optimization proceeds by maximizing the objective given in equation (\ref{eq:GRPO-obj}).

\subsubsection{Iterative RL with GRPO} During the reinforcement learning procedure, previously trained reward models may become inadequate for guiding the evolving policy model. To address this, we implement an iterative RL framework utilizing GRPO. As outlined in Algorithm \ref{alg:iter-grpo}, iterative GRPO involves generating updated training datasets for the reward model through samples produced by the current policy model. The reward model undergoes continual training using a replay strategy, whereby 10\% of prior data is incorporated. Subsequently, the policy model is re-designated as the reference model, and it is iteratively refined using the newly updated reward model.

\subsection{Training and Evaluating DeepSeekMath-RL}

Our RL experiments are based on DeepSeekMath-Instruct 7B, utilizing chain-of-thought-formatted questions from GSM8K and MATH extracted from the SFT dataset, totaling approximately 144K instances. We omit non-mathematical SFT questions in order to assess the effect of RL in scenarios where benchmark-related data is absent during RL. The reward model training set is curated in line with the methodology in \citep{wang2023math}. The initial reward model is fine-tuned from DeepSeekMath-Base 7B with a learning rate of 2e-5. For GRPO, we specify a learning rate of 1e-6 for the policy model and set the KL coefficient to 0.04. Each question is expanded into $64$ sampled completions. The sequence maximum length and batch size are both set to 1024. After each phase of exploration, the policy model receives a single update. DeepSeekMath-RL 7B is evaluated using the same benchmarks as DeepSeekMath-Instruct 7B. Within DeepSeekMath-RL 7B, tasks from GSM8K and MATH featuring chain-of-thought reasoning are considered in-domain, whereas the remaining benchmarks are classified as out-of-domain.

Table \ref{tab:sft_rl_math} presents a comparison of open- and closed-source models employing both chain-of-thought and tool-augmented reasoning on English and Chinese evaluation benchmarks. Our observations reveal that: 1) When employing chain-of-thought reasoning, DeepSeekMath-RL 7B achieves accuracies of 88.2\% on GSM8K and 51.7\% on MATH. These results exceed those of every open-source model ranging from 7B to 70B parameters, as well as most closed-source counterparts. 2) Notably, DeepSeekMath-RL 7B is solely trained on chain-of-thought-formatted instruction tuning data from GSM8K and MATH, initializing from DeepSeekMath-Instruct 7B. Even with this narrow training focus, DeepSeekMath-RL 7B outperforms DeepSeekMath-Instruct 7B in all assessed areas, underscoring the advantages introduced by reinforcement learning.

\section{Discussion}
In the following, we elaborate on the insights gained from both pre-training and reinforcement learning experiments.

\subsection{Lessons Learnt in Pre-Training}

We begin by sharing insights related to the pre-training process. Unless noted otherwise, training follows the protocols detailed in Section \ref{sec:quality-policy}. Throughout this section, references to the DeepSeekMath Corpus pertain to the 89B-token dataset constructed in the second iteration of data gathering.

\subsubsection{Code Training Benefits Mathematical Reasoning}

An often-discussed but yet to be substantiated conjecture is that training on code enhances reasoning skills. Our investigation contributes partial empirical evidence for this claim in the mathematical context: code training enhances models’ capabilities in mathematical reasoning tasks, with or without external tool assistance.

To assess the impact of code training on mathematical reasoning, we explored both two-stage and single-stage training paradigms:

\textbf{Two-Stage Training}

\begin{itemize}[topsep=0pt]
    \item \textbf{Code Training for 400B Tokens $\rightarrow$ Math Training for 150B Tokens}: DeepSeek-LLM 1.3B undergoes pretraining on 400B code tokens, which is subsequently followed by 150B tokens of math data;
    \item \textbf{General Training for 400B Tokens $\rightarrow$ Math Training for 150B Tokens}: For comparison, we also employ general tokens—sampled from DeepSeek-AI’s extensive general corpus—instead of code tokens for the initial training phase, in order to evaluate whether code-based pretraining offers any particular advantage in fostering mathematical reasoning compared to general-domain pretraining.
\end{itemize}

\textbf{One-Stage Training}

\begin{itemize}[topsep=0pt]
    \item \textbf{Math Training for 150B Tokens}: DeepSeek-LLM 1.3B is directly trained with 150B math tokens;
    \item \textbf{Training on a mixture of 400B Code Tokens and 150B Math Tokens}: Given that subsequent math training after code pretraining impairs coding abilities, we explore an alternative: jointly training the model on a combined dataset of code and math tokens to determine whether this approach can foster both mathematical reasoning and coding performance while counteracting catastrophic forgetting.
\end{itemize}

\paragraph{Results}
Table \ref{tab:code-math} and Table \ref{tab:code-math-general-eval} report the downstream outcomes obtained across these training paradigms.

In both two-stage and single-stage paradigms, code pretraining markedly promotes the model’s proficiency in program-augmented mathematical reasoning. Specifically, Table \ref{tab:code-math} illustrates that, under the two-stage pipeline, commencing with code tokens substantially strengthens performance on GSM8K and MATH tasks employing Python solutions. Additional math-specific training at the second stage further enhances this capability. For the one-stage design, integrating code and math tokens simultaneously in the training corpus reduces the severity of catastrophic forgetting observed in sequential training, and yields improvements in both code-based tasks (Table \ref{tab:code-math-general-eval}) and program-assisted mathematical problem solving (Table \ref{tab:code-math}).

Furthermore, code pretraining confers gains in mathematical reasoning even without explicit tool use. Within the two-stage configuration, preliminary code training yields modest but noticeable improvements and facilitates more effective math-specific fine-tuning, culminating in optimal results. However, blending code and math tokens for one-stage training appears to hinder mathematical reasoning without external tools. A plausible explanation is that the capacity constraints of DeepSeek-LLM 1.3B may restrict the model’s ability to concurrently internalize extensive code and math data.

\subsubsection{ArXiv Papers Appear to Offer Little Benefit for Mathematical Reasoning}

ArXiv papers have frequently been incorporated into math pre-training datasets \citep{minerva,gpt-f,llemma,mathpile}. Nevertheless, there remains a paucity of comprehensive studies investigating their direct influence on mathematical reasoning skills. Interestingly, our empirical results suggest that arXiv papers do not enhance mathematical reasoning capabilities as might be anticipated. Our investigation includes models at different scales—namely DeepSeek-LLM 1.3B and DeepSeek-Coder-Base-v1.5 7B \citep{deepseek-coder}—trained on arXiv datasets processed using distinct methodologies:

\begin{itemize}[topsep=0pt]
    \item \textbf{MathPile} \citep{mathpile}: a dataset containing 8.9 billion tokens, curated through cleaning and filtering based on heuristics, of which scientific arXiv articles constitute over 85\%;
    \item \textbf{ArXiv-RedPajama} \citep{redpajama}: a 28.0 billion-token corpus made up of all arXiv LaTeX files with ancillary content such as preambles, comments, macros, and bibliographies excised.
\end{itemize}

For our evaluations, DeepSeek-LLM 1.3B was trained on 150B tokens, while DeepSeek-Coder-Base-v1.5 7B was trained on 40B tokens from each of the arXiv corpora. The evidence points to arXiv documents being unhelpful for developing mathematical reasoning: when the models are exclusively trained on arXiv content, there is either no discernible improvement or an outright decline in their results across a spectrum of mathematical assessment tools employed here. The benchmarks include problem-solving datasets like GSM8K and MATH (Table \ref{tab:arxiv-cot}), standardized multiple-choice STEM questions such as MMLU-STEM (Table \ref{tab:arxiv-cot}), and formal proof environments like miniF2F (Table \ref{tab:arxiv_atp}).

Yet, it is important to acknowledge several caveats in our findings. Specifically, this analysis does not address:

\begin{itemize}[topsep=0pt]
    \item How arXiv-sourced content might affect mathematical tasks absent from this work, for example, the informal exposition of formal results or proofs;
    \item The influence of arXiv data when intermingled with other forms of pre-training material;
    \item Whether scaling up model size could reveal benefits that are not present at the evaluated sizes.
\end{itemize}

Consequently, these unanswered questions merit further inquiry in future work.

\subsection{Insights of Reinforcement Learning}
\subsubsection{Towards to a Unified Paradigm} In this section, we propose a unified analytical framework encompassing various training methodologies, such as SFT, RFT, DPO, PPO, and GRPO, and present empirical studies probing the variables of this framework. Typically, the parameter gradient $\theta$ for any given training approach can be expressed as:

\begin{equation}
    \nabla_{\theta}\mathcal{J}_{\textcolor{red}{\mathcal{A}}}(\theta) = \mathbb{E}[\underbrace{(q,o) \sim \textcolor{red}{\mathcal{D}}}_{Data \ Source}]\left( \frac{1}{|o|} \sum_{t=1}^{|o|}  \underbrace{GC_{{\mathcal{A}}}(q, o, t, \textcolor{red}{\pi_{{rf}}})}_{Gradient \ Coefficient}  \nabla_{\theta}\log \pi_{\theta}(o_t | q, o_{<t})\right).
\label{eq:objective}
\end{equation}

The framework is characterized by three principal elements: 1) the \textit{Data Source} $\mathcal{D}$, specifying the set of examples for training; 2) the \textit{Reward Function} $\pi_{{rf}}$, which provides the feedback signal essential for learning; and 3) the \textit{Algorithm} $\mathcal{A}$, which utilizes both the data and reward signal to produce the gradient coefficient $GC$, thereby modulating the degree of penalization or reward applied during training. Within this structure, several important approaches can be outlined:

\begin{itemize}
    \item  \textbf{Supervised Fine-tuning (SFT)}: In SFT, a pretrained model undergoes further training on data curated by human annotators.
    \item \textbf{Rejection Sampling Fine-tuning (RFT)}: RFT takes the SFT-tuned model and continues training it using a subset of its generated answers, where outputs are retained based on their correctness relative to SFT prompts.
    \item \textbf{Direct Preference Optimization (DPO)}: DPO enhances the SFT-trained model by optimizing on augmented samples, with a pairwise loss structure designed to reflect human preferences.
    \item \textbf{Online Rejection Sampling Fine-tuning (Online RFT)}: Unlike standard RFT, Online RFT iteratively updates the policy model—initialized from the SFT model—by fine-tuning it on outputs sampled dynamically from the model itself.
    \item \textbf{PPO/GRPO}: For PPO/GRPO, the SFT model serves as the starting point for the policy, which is subsequently improved through reinforcement using outputs generated by the online policy.
\end{itemize}

The components of these methods are outlined in Table \ref{tab:data_policy}. For a comprehensive explanation of the derivation steps, see Appendix \ref{app:analysis-rl}.

\paragraph{Observation about Data Source} We categorize data sources into two types: online sampling and offline sampling. Online sampling refers to collecting training data generated through exploration by the current training policy model in real time, whereas offline sampling implies using training data derived from samples produced by the original SFT model. Both RFT and DPO utilize the offline sampling approach, while Online RFT and GRPO adopt the online strategy.

As depicted in Figure \ref{fig:rl-analysis}, our results indicate that Online RFT consistently achieves higher performance than RFT across two evaluation benchmarks. Notably, in the initial training phase, Online RFT performs similarly to RFT, but as training progresses, Online RFT secures a clear performance lead, illustrating the benefits of online training. This pattern is expected: early in training, the actor and SFT model remain largely similar, so the samples from the actor exhibit only minor divergence. At later stages, however, samples from the evolving actor display more substantial variations, and the use of up-to-date online samples becomes increasingly advantageous.

\paragraph{Observation about Gradient Coefficient} In these algorithms, the transformation of the reward signal into a gradient coefficient governs how model parameters are updated. Our experiments distinguish between two reward mechanisms: `Rule' and `Model'. The `Rule' approach evaluates the quality of responses using explicit answer correctness, whereas the `Model' approach involves training a separate reward model (whose training data is derived from rule-based judgments) to assign scores to responses. Equations \ref{eq:GC-RFT} and \ref{eq:GC-GRPO} reveal a critical distinction: GRPO uniquely modulates its gradient coefficient in accordance with the reward model’s outputs, thereby applying varied strengths of reinforcement or penalization contingent on the score. In comparison, Online RFT does not incorporate this adaptive adjustment; it provides a consistent level of positive reinforcement for all correct answers and omits penalties for incorrect responses.

As illustrated in Figure \ref{fig:rl-analysis}, GRPO achieves better results than online RFT, underlining the effectiveness of modifying the coefficients for positive and negative gradients. Additionally, GRPO+PS outperforms GRPO+OS, which underscores the advantage of applying more granular, step-sensitive gradient coefficients. Our investigation also extends to iterative RL; specifically, we perform two iterations in our study. According to Figure \ref{fig:iter-rl}, iterative RL notably boosts performance, with the first iteration providing the most substantial gains.

\subsubsection{Why RL Works?} This work applies reinforcement learning on a subset of instruction tuning data and observes marked performance gains over models trained solely with instruction tuning. To better understand the mechanisms behind this improvement, we analyze Pass@K and Maj@K accuracies for both the Instruct and RL models across two evaluation benchmarks. Figure \ref{fig:maj-pass} reveals that RL drives an increase in Maj@K but leaves Pass@K largely unchanged. This suggests that RL does not directly enhance base capabilities but instead increases the likelihood that a correct answer appears among the top $K$ responses—\textbf{thus, improvements seem to result from amplifying the prominence of correct outputs within TopK rather than strengthening core competencies.} Similarly, prior work \citep{wang2023making} reports a \textbf{misalignment problem} with SFT models on reasoning tasks and demonstrates that preference alignment methods \citep{yuan2023rrhf,song2023preference,wang2023making} can significantly improve the reasoning abilities of such models.

\subsubsection{How to Achieve More Effective RL?}
\label{sec:effec_rl}
Our results demonstrate that RL is highly effective for mathematical reasoning tasks, and we also propose a unified perspective for interpreting different representative training approaches. In this framework, each method may be seen as either an instance of RL or a simplified variant. Equation \ref{eq:objective} highlights three principal elements: Data Source, Algorithm, and Reward Function. We suggest future research avenues for each component.

\paragraph{Data Source} The data source underpins all forms of training. For RL, the data source refers specifically to the set of unlabeled queries, coupled with responses sampled from the policy model. In this work, we limit ourselves to using questions from the instruction tuning phase, with response candidates generated via basic nucleus sampling. This approach may partly explain why our RL pipeline mainly leads to improvements in Maj@K. Moving forward, we intend to extend our RL pipeline to questions drawn from distributions not seen during instruction tuning and to adopt \textbf{advanced sampling (decoding) strategies}, including methods informed by tree search \citep{yao2023tree}. Furthermore, the integration of \textbf{efficient inference techniques} \citep{xia-etal-2023-speculative,leviathan2023fast,kwon2023efficient,xia2024unlocking}, which shape how effectively the policy model explores the data, will also be of critical significance.

\paragraph{Algorithms} Algorithms utilize both the observed data and the received reward signal to calculate the gradient coefficient used for updating the model parameters. Referring to Equation \ref{eq:objective}, it can be observed that current approaches largely place complete \textbf{TRUST} in the reward function signal when deciding whether to promote or suppress the conditional probability assigned to a token. Nonetheless, such unwavering trust is problematic, as the reliability of reward signals cannot always be guaranteed, particularly for highly challenging tasks. For instance, the PRM800K datasets \citep{lightman2023let}, despite thorough annotation by experienced annotators, still exhibit roughly 20\% erroneous labels\footnote{\url{https://github.com/openai/prm800k/issues/12\#issuecomment-1728491852}}. Given these challenges, we focus on investigating reinforcement learning algorithms that maintain robustness in the presence of noisy reward signals. We propose that the adoption of \textbf{WEAK-TO-STRONG} alignment techniques \citep{burns2023weak} may introduce profound shifts in the development of learning algorithms.

\paragraph{Reward Function} The reward function serves as the primary provider of learning signals during training. In the context of RL, the reward signal typically originates from a neural reward model. We identify three key areas of investigation for reward models: 1) \textbf{Strategies for improving the reward model's generalization capabilities.} Ensuring that the reward model generalizes well to unseen question types and novel decoding outputs is critical; failure to do so risks RL merely reinforcing existing model behaviors, rather than facilitating substantive improvement; 2) \textbf{Methods to quantify and incorporate reward model uncertainty.} Capturing model uncertainty can potentially facilitate the interaction between a less robust reward model and weak-to-strong learning approaches; 3) \textbf{Approaches to the efficient construction of high-quality process reward models} capable of delivering precise, step-level training signals for the reasoning process \citep{lightman2023let,wang2023math}.

\section{Conclusion, Limitation, and Future Work}

We introduce DeepSeekMath, which sets a new standard for open-source models on the challenging MATH benchmark and narrows the gap with proprietary systems. DeepSeekMath~builds upon DeepSeek-Coder-v1.5 7B as its foundation and is continually trained on 500B tokens, including a crucial subset of 120B math-specific tokens extracted from Common Crawl. Our comprehensive ablation analysis demonstrates that web pages represent a valuable source for high-quality mathematical data, whereas arXiv did not contribute as significantly as anticipated. We propose Group Relative Policy Optimization (GRPO), an alternative to Proximal Policy Optimization (PPO), which enables marked improvements in mathematical reasoning with reduced memory requirements. Experimental evidence indicates that GRPO is beneficial even when DeepSeekMath-Instruct 7B attains strong results on standard benchmarks. Additionally, we offer a consolidated framework for interpreting various methodologies and discuss several promising avenues for more efficient reinforcement learning strategies.

Despite DeepSeekMath's notable advancement on quantitative reasoning assessments, its proficiency in geometry and theorem-proving remains lower than that of closed-source systems. As an example, in preliminary evaluations, the model failed on questions involving geometric entities like triangles and ellipses, pointing to a possible data selection bias during pre-training and fine-tuning stages. Furthermore, due to model size constraints, DeepSeekMath~does not match GPT-4’s capabilities in few-shot learning settings; while GPT-4’s performance benefits from few-shot prompts, DeepSeekMath~performs similarly in both zero-shot and few-shot evaluations. Going forward, we aim to refine our curated data selection workflow to yield an even more robust pre-training dataset. We will also pursue the potential pathways outlined in Section \ref{sec:effec_rl} to realize more effective reinforcement learning approaches for large language models.

\end{document}